\documentclass[12pt,halfparskip]{scrreprt}

\usepackage{pystructure}

\usepackage{svn} % For handling of SVN keywords
\SVN $LastChangedDate: 2008-01-25 14:34:07 +0100 (Fr, 25 Jan 2008) $

\begin{document}

\begin{titlepage}
\begin{center}
\thispagestyle{empty}

\begin{figure}[h]
 \centering
 \vspace{0,5cm}
 \includegraphics[width=\textwidth]{img/hsr_logo}
\end{figure}

\vspace{1cm}
{\Large \bfseries HSR -- University of Applied Sciences Rapperswil}

\vspace{0,5cm}
{\Large \bfseries Institute for Software}

\vspace{2cm}
{\Huge \bfseries PyStructure -- Automated Structure and Dependency Analysis of Python Code}
\vspace{2cm}

{\Large \bfseries Bachelor Thesis: Spring 2008}

\vspace{0,5cm}
\SVNDate


\vspace{1cm}
Reto Schüttel \\ \url{reto@schuettel.ch}

\vspace{0,5cm}
Robin Stocker \\ \url{robin@nibor.org}

\vspace{0,5cm}
Supervised by Prof. Peter Sommerlad

% FIXME external partner

\vspace{1cm}
\url{http://pystructure.ifs.hsr.ch/}

\end{center}
\end{titlepage}


\pagenumbering{roman}
\pagestyle{plain}

%%%%%%%%%%%%%%%%%%%%%%%%%%%%%%%%%%%%%%%%%%%%%%%%%%%%%%%%%%%%%%%%%%%%%%%
\chapter*{Abstract}


%%%%%%%%%%%%%%%%%%%%%%%%%%%%%%%%%%%%%%%%%%%%%%%%%%%%%%%%%%%%%%%%%%%%%%%
\chapter*{Management Summary}

\section*{Motivation}

\section*{Goal}

\section*{Results}

\section*{Outlook}


\newpage

\tableofcontents

\newpage
\pagenumbering{arabic}
\pagestyle{scrheadings}

%%%%%%%%%%%%%%%%%%%%%%%%%%%%%%%%%%%%%%%%%%%%%%%%%%%%%%%%%%%%%%%%%%%%%%%
\chapter{Introduction}

\section{Structural Analysis}

Structural Analysis

%%%%%%%%%%%%%%%%%%%%%%%%%%%%%%%%%%%%%%%%%%%%%%%%%%%%%%%%%%%%%%%%%%%%%%%
\chapter{Type Inference}

\section{Introduction}

Idea \& Concept \\
goal, ziel (wo wollen wir hin)\\
(Goal engine, why \& how) \\
(link to other possible solutions)
Correctness \& Usefulness \\
Process Overview \\

\section{Parsing the Source Code}

%{ Korrektur&Feedback Robin 
Before any analysis can be done on the application its source code has to be brought into form which is easier to process for programs. It's quite common to express an application's source code as a tree, where each interior node represents a programming language construct and the children of that node represent meaningful components of the construct. For example the simple code \code{foo("bar")} is a function node with the parameter as a string node attached to it. This structure is usually called Abstract Syntax Tree (AST).

Instead of writing a new parser PyStructure uses a modified version of the Jython\footnote{Jython is a Python implementation written in Java} parser. The modifications were done by the PyDev(TODO footnote) project to be able to interpret Python 2.5 source code, the normal Jython parser can only parse source code written for Python 2.2. In the long term it would be advantageous to use the unmodified Jython parser, but before 2.5 syntax is supported this is not really an option.
%}

\section{Creating the Model}
From AST to Model \\
(StructureVisitor, DefinitionVisitor)\\
Connecting Definitions and Uses\\
(What's already known, what not (attributes...))\\


\section{Inferencing Types}
Goal Engine \\
 (description of the goal engine)\\
 (advantages/disadvantages)\\
Evaluation\\
 (goals \& evaluators explained)\\
 (examples)\\


\section{Challenges/Problems}
Container Element Types \\
Call Context \& Instance Context \\
Built-in Types \& Functions\\
Inheritance\\
 (C3 MRO)\\
Recursion\\

\section{Caching}
 Techniques, Performance, advantages/disadvantages, problems
 Result, problems, performance, memory requirements (ganzer ti)
\section{Improvements During the Development}
 (DLTK \& goalengine refactoring) \\
 (Avoid casting (result in Goal))\\
 (Goaltype -> Evaluator Map)\\
Other Type Inference Solutions\\

%%%%%%%%%%%%%%%%%%%%%%%%%%%%%%%%%%%%%%%%%%%%%%%%%%%%%%%%%%%%%%%%%%%%%%%
\chapter{Structural Analysis}
How to interpret the information provided by the TI \\
What kind of relations \\
What kind of elements (Module, class, etc. mapping to s101g) \\
How does the output look like \\
Parsing a single expression vs inference all expressions (demand driven vs everything)\\
 what are `all expressions`\\

%%%%%%%%%%%%%%%%%%%%%%%%%%%%%%%%%%%%%%%%%%%%%%%%%%%%%%%%%%%%%%%%%%%%%%%
\chapter{Results / Practical examples}
Structure \\
types (type annotator)\\
Key figures\\



%% What about tthis chapter? %%%%%%%%%%%%%%%%%%%%%%%%%%%%%%%%%%%%%%%%%%%%%%%%%%%%%%%%%%%%%%%%%%%%%%%
\chapter{Structural Improvements}
\section{Size Over Time}
\section{Overall Structure}
\section{Type Inference Engine}
\section{Quality Assurance Infrastructure}
\subsection{Continuous Integration}
\subsection{Automatic Unit Testing}
\subsection{Checkstyle}
\subsection{Findbugs}
\section{Results}

\chapter{Continuous Integration \& Testing}
 Ways to test TI
 Framework, markers
  (examples, mro marker)
 Continuous Integration
  (lava lamp)

\chapter{Outlook}
 Possible follow-up projects
 Possibilities.. etc.

\chapter{Project Management}
 Milestones \& Roadmap \\
 (Problems, important decisions) \\
 (QA Techniques) \\
 (Used license) \\
 Effort Diagram \\
 Our Experiences \\
 Acknowledgements 

\section{Outlook}
\subsection{PEPTIC in General}
\subsection{Type Inference}
\subsubsection{Common IDE Features}
\subsubsection{Type Checking}
\subsubsection{Structural Analysis}
\subsubsection{Unused Code}
\subsubsection{Compiler Optimisations}

\clearpage

\section{Effort Diagram}

\clearpage

\section{Thanks}

\begin{itemize}
	\item to Peter Sommerlad and Thomas Corbat for their support during the project
	\item to Christian Bachmann for the life-saving coffee machine in our study room. 
	\item NOT to HSR for their banning of the above mentioned life-saving coffee machine out of the study room.	
\end{itemize}


%%%%%%%%%%%%%%%%%%%%%%%%%%%%%%%%%%%%%%%%%%%%%%%%%%%%%%%%%%%%%%%%%%%%%%%

\listoffigures

\bibliographystyle{alphadin}
\bibliography{bibliography}


\end{document}
